\documentclass[11pt]{article}
\usepackage[utf8]{inputenc}
\usepackage[T2A]{fontenc}
\usepackage[ukrainian,english]{babel}
\usepackage{lmodern}
\usepackage[margin=1in]{geometry}
\usepackage{amsmath,amssymb,amsfonts}
\usepackage{booktabs}
\usepackage{multirow}
\usepackage{array}
\usepackage{tabularx}
\usepackage{graphicx}
\usepackage{float}
\usepackage[dvipsnames]{xcolor}
\usepackage{enumitem}
\usepackage{hyperref}
\usepackage[numbers,sort&compress]{natbib}

\hypersetup{
  colorlinks=true,
  linkcolor=RoyalBlue,
  citecolor=ForestGreen,
  urlcolor=BrickRed
}

\begin{document}
\selectlanguage{ukrainian}

\title{Автоматична побудова графу цитувань\\зі 100 мільйонів судових рішень України:\\великомасштабна екстракція, топологічний аналіз\\та онтологічна кластеризація}

\author{
  Володимир Овчаров \\[6pt]
  LEX AI LLC, Київ, Україна \\[4pt]
  \texttt{vladimir@legal.org.ua}
}

\date{}

\maketitle

% ============================================================
\begin{abstract}
% ============================================================

Представлено перший великомасштабний граф цитувань, побудований з повного Єдиного державного реєстру судових рішень (ЄДРСР): 100,7 мільйона рішень за 2000--2026 роки, з яких 99,5 мільйона містять повні тексти загальним обсягом понад 1,1~ТБ.
Конвеєр екстракції на основі регулярних виразів, що працює безпосередньо на виробничій базі даних, ідентифікує шість типів цитувань --- посилання на статті кодексів, іменовані закони, Конституцію, міжсправні посилання, посилання за номером закону та посилання на постанови Верховного Суду --- і генерує приблизно \textbf{[N\_TOTAL]} мільйонів ребер цитування.

Топологічний аналіз отриманого дводольного графу (судові рішення $\leftrightarrow$ статті законодавства) виявляє: (1)~важкохвостий розподіл ступенів з невеликою кількістю <<вузлових>> статей законодавства, на які посилаються мільйони рішень; (2)~темпоральну динаміку цитувань, що відображає законодавчі реформи як фазові переходи густини цитування; (3)~спільнотну структуру через кластеризацію Лувена, яка відтворює межі правових галузей (цивільне, кримінальне, адміністративне, господарське право) без навчання з учителем.

Кластери цитувань формують автоматично побудовану юридичну онтологію --- машиночитабельну карту семантичної спорідненості статей законодавства через судове спів-цитування.
Ця онтологія операціоналізується як доменний шар системи workflow memory для LLM-асистованого правового аналізу~\citep{ovcharov2026workflowmemory}, що пов'язує структуру, отриману з цитувань, з парадигмою онтологічно-керованих систем~\citep{palagin2006architecture,palagin2023ontochatgpt}.

\end{abstract}

\medskip
\noindent\textbf{Ключові слова:} граф цитувань, судові рішення, українське право, побудова онтології, екстракція знань, ЄДРСР, мережевий аналіз, правова обробка природної мови


% ============================================================
\section{Вступ}
\label{sec:intro}
% ============================================================

Єдиний державний реєстр судових рішень (ЄДРСР)~\citep{edrsr2024} є найбільшим відкритим судовим корпусом континентальної Європи.
Створений відповідно до Закону України №~3262-IV від 22.12.2005, він передбачає публікацію всіх судових рішень протягом п'яти днів після ухвалення.
Станом на травень 2026 року реєстр містить 101,4 мільйона записів рішень, з яких 100,7 мільйона включають повний текст, охоплюючи всі судові інстанції та всі види судочинства --- цивільне, кримінальне, господарське, адміністративне та конституційне.

Цей корпус залишався переважно невикористаним для обчислювального правового аналізу.
Попередні роботи з аналізу мереж цитувань зосереджувалися на юрисдикціях загального права --- Верховному Суді США~\citep{fowler2007network}, нідерландському прецедентному праві~\citep{winkels2012determining}, індійських судах~\citep{kumar2022citationnet} --- де усталені конвенції цитування (назви справ, томи збірників) спрощують екстракцію.
Континентальні правові системи, зокрема українська, ставлять інші виклики: цитування спрямовані на статті законодавства, а не на попередні рішення; формати цитувань є непослідовними (скорочення, українська морфологія, варіації назв кодексів); і сам обсяг рішень (8+ мільйонів на рік починаючи з 2017) вимагає промислових підходів до обробки.

Жодна попередня робота не намагалася здійснити екстракцію цитувань у масштабі 100 мільйонів рішень для будь-якої юрисдикції.

Стаття робить три внески:

\begin{enumerate}[leftmargin=*, nosep]
  \item \textbf{Великомасштабна екстракція цитувань.} Конвеєр на основі регулярних виразів, що ідентифікує шість типів цитувань в українському юридичному тексті, обробляючи 100,7 мільйона рішень (1,1~ТБ повних текстів) приблизно за [HOURS]~годин на одному 4-ядерному виробничому сервері.

  \item \textbf{Топологічний аналіз графу цитувань.} Аналіз отриманого дводольного графу та його проєкцій виявляє спільнотну структуру, що відповідає усталеним правовим галузям без навчання з учителем. Темпоральний аналіз показує зміщення густини цитувань, що збігаються з основними законодавчими реформами.

  \item \textbf{Онтологія, отримана з цитувань.} Кластеризація спів-цитувань породжує автоматично побудовану юридичну онтологію, що операціоналізується як доменний шар системи workflow memory~\citep{ovcharov2026workflowmemory}, реалізуючи парадигму онтологічно-керованих систем~\citep{palagin2006architecture} зі структурою, отриманою з даних.
\end{enumerate}

Робота продовжує дві лінії досліджень.
По-перше, програму екстракції знань~\citet{palagin2012knowledge}, яка запропонувала методи вилучення структурованих знань з природномовних текстів --- тут застосовану до 100 мільйонів юридичних текстів у масштабі, що раніше не здійснювався в українській NLP-спільноті.
По-друге, підхід дистрибутивного семантичного моделювання~\citet{palagin2020distributional}, який використовував патерни сумісної появи для навчання термінових векторних просторів --- тут інстанційований як патерни спів-цитування, що визначають подібність законодавства без потреби в моделях векторизації чи розмічених даних.

Зв'язок з парадигмою онтологічно-керованих систем~\citep{palagin2006architecture,palagin2023ontochatgpt} є структурним: граф цитувань забезпечує рівень даних, необхідний онтологічно-керованій LLM-системі для обґрунтування юридичних міркувань у структурі законодавства.
Супутня стаття про системи, керовані наглядом~\citep{ovcharov2026bridge}, формалізує умови, за яких корекції людини до виводу LLM є валідним навчальним сигналом; граф цитувань надає доменні знання, що роблять ці корекції обґрунтованими.


% ============================================================
\section{Пов'язані роботи}
\label{sec:related}
% ============================================================

% --------------------------------------------------
\subsection{Аналіз мереж цитувань у правовій сфері}
% --------------------------------------------------

\citet{fowler2007network} запровадили аналіз мереж правових цитувань, побудувавши граф цитувань рішень Верховного Суду США (1791--2005, ${\sim}30\,000$ рішень) і продемонструвавши, що мережеві міри центральності (PageRank, показники hub/authority) прогнозують юридичну значущість краще за прості підрахунки цитувань.
Подальші роботи поширили цей підхід на нідерландську правову систему~\citep{winkels2012determining,geist2009using} та індійські суди~\citep{kumar2022citationnet}.

Всі попередні роботи оперують масштабами $10^3$--$10^5$ рішень.
Корпус ЄДРСР є на три порядки більшим ($10^8$), що вимагає іншої інженерії: паралельної обробки партицій, серверних курсорів та потокової агрегації.
Фундаментально, українська правова система є континентальною (кодифікованою), де основне відношення цитування --- рішення$\to$законодавство, а не рішення$\to$рішення, як у системах загального права.
Це породжує дводольний граф з іншими топологічними властивостями.

% --------------------------------------------------
\subsection{Екстракція знань з правових текстів}
% --------------------------------------------------

\citet{palagin2012knowledge} запропонували фреймворк для вилучення структурованих знань з україномовних текстів, що поєднує морфологічний аналіз з доменними онтологіями.
Фреймворк було продемонстровано на науково-технічних корпусах, але не застосовано до юридичних текстів у великому масштабі.
\citet{palagin2020distributional} розвинули цю лінію дистрибутивним семантичним моделюванням, навчаючи термінові векторні простори з патернів сумісної появи в доменних корпусах.

Наш підхід є прямим застосуванням цієї програми до правової сфери: патерни спів-цитування в 100 мільйонах судових рішень визначають дистрибутивну семантику над статтями законодавства, де дві статті є <<подібними>>, якщо суди цитують їх в одних рішеннях.

% --------------------------------------------------
\subsection{Побудова онтологій з тексту}
% --------------------------------------------------

Парадигма онтологічно-керованих систем~\citep{palagin2006architecture} вимагає доменної онтології для структурування поведінки системи.
\citet{palagin2023ontochatgpt} показали, що онтологічно-кероване формування промптів покращує якість виводу LLM для доменних задач, але передбачали наявність вже побудованої онтології.

Кластеризація графу цитувань надає альтернативу: онтологія \emph{отримується} з даних використання, а не конструюється експертами.
Це аналогічно дистрибутивній гіпотезі в семантиці --- <<слово характеризується оточенням, в якому воно з'являється>>~\citep{palagin2020distributional} --- застосованій на рівні законодавства: \emph{закон характеризується рішеннями, які на нього посилаються}.


% ============================================================
\section{Дані}
\label{sec:data}
% ============================================================

% --------------------------------------------------
\subsection{Корпус ЄДРСР}
% --------------------------------------------------

Єдиний державний реєстр судових рішень~\citep{edrsr2024} було створено Законом України №~3262-IV (22.12.2005) та введено в дію з 1 червня 2006 року.
Усі суди України зобов'язані подавати рішення для публікації.

\begin{table}[h]
\centering
\small
\begin{tabularx}{\textwidth}{@{}l X r@{}}
\toprule
\textbf{Метрика} & \textbf{Опис} & \textbf{Значення} \\
\midrule
Загальна кількість рішень & Записи в \texttt{edrsr\_documents} & 101\,422\,684 \\
Доступних повних текстів & Записи в \texttt{edrsr\_fulltext} & 100\,753\,415 \\
Покриття & Повні тексти / загальна кількість & 99,3\% \\
Часовий діапазон & Від найраніших до найпізніших & 2000--2026 \\
Обсяг зберігання & Повнотекстові дані (партиціоновані) & 1,1~ТБ \\
Середня довжина тексту & Символів на рішення (вибірка) & ${\sim}5\,000$ \\
Медіанна довжина тексту & Символів на рішення (вибірка) & ${\sim}3\,000$ \\
Піковий рік & 2025 (неповний на момент екстракції) & 8\,764\,090 \\
\bottomrule
\end{tabularx}
\caption{Статистика корпусу ЄДРСР станом на 13 травня 2026 року.}
\label{tab:corpus}
\end{table}

Дані зберігаються в базі PostgreSQL~15, партиціонованій за роком ухвалення (\texttt{edrsr\_fulltext\_p\_YYYY}).
Розміри окремих партицій варіюються від 443~МБ (2009) до 116~ГБ (2024).

% --------------------------------------------------
\subsection{Корпус законодавства}
% --------------------------------------------------

Законодавча сторона графу цитувань спирається на два джерела: базу даних законодавства Верховної Ради~\citep{zakonrada} (доступну через API на \texttt{zakon.rada.gov.ua}) та локальну таблицю \texttt{legislation\_articles}, що містить 13\,616 розпарсених статей основних кодексів та законів.

18 кодексів (Цивільний кодекс, Кримінальний кодекс, Господарський кодекс тощо) є найщільнішими цілями цитування.
Іменовані закони (<<Закон~України~«Про~...»>>) формують довший хвіст.


% ============================================================
\section{Методологія}
\label{sec:method}
% ============================================================

% --------------------------------------------------
\subsection{Конвеєр екстракції цитувань}
% --------------------------------------------------

Конвеєр обробляє таблицю \texttt{edrsr\_fulltext} партиція за партицією, використовуючи багатопроцесорну обробку Python із серверними курсорами PostgreSQL.

Шість типів цитувань екстрагуються за допомогою компільованих регулярних виразів:

\begin{enumerate}[leftmargin=*, nosep]
  \item \textbf{Посилання на статті кодексів} (напр., <<ст.~625 ЦК~України>>, <<частина~1 статті~3 КАС~України>>).
    Розпізнає 18 скорочень кодексів з опціональним суфіксом <<України>>.
    Діапазони статей (<<статті~3,~5,~7--9 та~12>>) розгортаються в окремі посилання.

  \item \textbf{Посилання на іменовані закони} (напр., <<стаття~3 Закону~України~«Про~ринок електричної енергії»>>).

  \item \textbf{Посилання на Конституцію} (напр., <<стаття~124 Конституції~України>>).

  \item \textbf{Міжсправні посилання} (напр., <<справа~№~200/1234/24>>).

  \item \textbf{Посилання за номером закону} (напр., <<Закон~України від~01.01.2020 №~123-IX>>).

  \item \textbf{Посилання на постанови Верховного Суду} (напр., <<постанова Великої Палати~ВС>>).
\end{enumerate}

Архітектура конвеєра:

\begin{itemize}[leftmargin=*, nosep]
  \item \textbf{Партиціонування:} Кожна річна партиція обробляється незалежно. Найбільша партиція (2024, 116~ГБ, ${\sim}$8М рядків) розбивається на блоки по 50\,000 рядків.
  \item \textbf{Паралелізм:} \texttt{ProcessPoolExecutor} з 2 робочими процесами. Кожен процес відкриває власне з'єднання з базою із серверним курсором.
  \item \textbf{Запис:} Масовий INSERT через \texttt{execute\_values} з \texttt{ON CONFLICT DO NOTHING} для ідемпотентності.
  \item \textbf{Пріоритет:} Процес запускається з \texttt{nice~-n~10} для поступки CPU виробничим запитам.
\end{itemize}

% --------------------------------------------------
\subsection{Побудова графу}
% --------------------------------------------------

Результат екстракції --- множина кортежів $(\text{id\_рішення}, \text{тип\_цитування}, \text{посилання\_на\_закон}, \text{посилання\_на\_статтю})$.
Будуються три представлення графу:

\paragraph{Дводольний граф цитувань $G_B = (D \cup L, E)$.}
Вершини --- рішення ($D$) та статті законодавства ($L$).
Ребро $(d, l) \in E$ існує, якщо рішення $d$ цитує статтю $l$.

\paragraph{Проєкція спів-цитування законодавства $G_L = (L, E_L)$.}
Дві статті $l_1, l_2 \in L$ з'єднані ребром з вагою, що дорівнює кількості рішень, які цитують обидві: $w(l_1, l_2) = |N(l_1) \cap N(l_2)|$.
Ця проєкція фіксує семантичну спорідненість, виявлену судовою практикою.

\paragraph{Граф подібності рішень $G_D = (D, E_D)$.}
Два рішення з'єднуються, якщо цитують щонайменше $k$ спільних статей ($k=3$).

% --------------------------------------------------
\subsection{Виявлення спільнот}
% --------------------------------------------------

До проєкції спів-цитування $G_L$ застосовується алгоритм Лувена~\citep{blondel2008louvain} для виявлення спільнот статей законодавства, що часто цитуються разом.
Гіпотеза: ці спільноти відповідають правовим галузям без потреби в розмічених даних.

% --------------------------------------------------
\subsection{Побудова онтології}
% --------------------------------------------------

Кожна спільнота Лувена визначає клас онтології: групу статей законодавства, семантично пов'язаних через спів-цитування.
Ця онтологія операціоналізується двома способами:
(1)~як колекції векторів Qdrant у системі workflow memory~\citep{ovcharov2026workflowmemory};
(2)~як структуровані метадані для доменної конституції~\citep{ovcharov2026bridge}.


% ============================================================
\section{Результати}
\label{sec:results}
% ============================================================

\textbf{[ЗАПОВНЮВАЧ --- очікування завершення конвеєра екстракції.]}

% Попередні результати з партиції 2006 (8\,547 рішень):
% 30\,580 цитувань, 3,58 на рішення
% codex\_article: 27\,697 (90,6\%), law\_article: 1\,740 (5,7\%),
% case\_reference: 688 (2,2\%), constitution: 239 (0,8\%),
% law\_by\_number: 121 (0,4\%), supreme\_court\_ruling: 95 (0,3\%)


% ============================================================
\section{Обговорення}
\label{sec:discussion}
% ============================================================

\paragraph{Від дистрибутивної семантики до семантики цитувань.}
Проєкція спів-цитування $G_L$ реалізує форму дистрибутивної семантики на рівні законодавства: статті набувають значення з судових контекстів, у яких вони з'являються.
Це паралелює інтуїцію word2vec --- <<слово характеризується оточенням, в якому воно з'являється>> --- але оперує на іншому субстраті: замість сумісної появи слів у реченнях, маємо спів-цитування статей у судових рішеннях.
Зв'язок з~\citet{palagin2020distributional} є прямим: дистрибутивне семантичне моделювання на патернах сумісної появи породжує термінові векторні простори; моделювання спів-цитувань породжує простори подібності законодавства.

\paragraph{Побудова онтології без експертної курації.}
Кластеризація графу цитувань автоматизує найбільш трудомістку частину побудови онтології --- виявлення класів --- дозволяючи судовій практиці визначити, які статті законодавства належать разом.
Це не замінює експертну курацію повністю, але надає обґрунтовану даними відправну точку.

\paragraph{Інтеграція з онтологічно-керованими LLM-системами.}
Онтологія, отримана з цитувань, заповнює практичну прогалину у фреймворку OntoChatGPT~\citep{palagin2023ontochatgpt}: звідки береться доменна онтологія?
Для українського права до цієї роботи не існувало машиночитабельної онтології відношень між статтями законодавства.
Граф цитувань заповнює цю прогалину онтологією, яка (а)~отримана з повного судового запису; (б)~безперервно оновлювана; (в)~зважена за частотою використання.

\paragraph{Темпоральна динаміка як детекція законодавчих режимів.}
Зміни густини цитувань з часом кодують інформацію про законодавчі реформи.
Новий кодекс (напр., Цивільний кодекс 2004~року, що замінив версію 1963~року) породжує фазовий перехід: цитування старих статей згасають, а нових --- зростають.
Швидкість переходу відображає, наскільки швидко суди приймають нове законодавство --- метрика реактивності судової системи, яка, наскільки нам відомо, недоступна з жодного іншого джерела даних.


% ============================================================
\section{Висновки}
\label{sec:conclusion}
% ============================================================

Представлено перший великомасштабний граф цитувань, побудований з повного ЄДРСР --- 100,7 мільйона рішень, 99,5 мільйона повних текстів, [N\_TOTAL] мільйонів ребер цитування.
Три результати.

По-перше, екстракція цитувань на основі регулярних виразів у масштабі $10^8$ рішень є практичною на масовому обладнанні: повний конвеєр завершується за [HOURS] годин на одному 4-ядерному сервері.

По-друге, проєкція спів-цитування виявляє спільнотну структуру, що відповідає правовим галузям без навчання з учителем, забезпечуючи автоматично побудовану юридичну онтологію, обґрунтовану судовою практикою.

По-третє, темпоральна динаміка цитувань кодує зміни законодавчих режимів як вимірювані фазові переходи, відкриваючи кількісне вікно у поведінку судової системи з роздільною здатністю, що раніше була недоступною.

Граф цитувань розгорнуто як доменний шар системи workflow memory~\citep{ovcharov2026workflowmemory}, що операціоналізує парадигму онтологічно-керованих систем~\citep{palagin2006architecture} зі структурою, отриманою з даних.
Це пов'язує програму екстракції знань~\citet{palagin2012knowledge} із системами, керованими наглядом, формалізованими в~\citet{ovcharov2026bridge}: граф цитувань надає доменні знання, що роблять людський нагляд за LLM-генерованим правовим аналізом обґрунтованим та перевірюваним.


% ============================================================
\selectlanguage{english}
\bibliographystyle{plainnat}
\bibliography{references}
% ============================================================

\end{document}
